%% ===================================================================
\section{Summary and outlook}\label{sec:summary}
%% ===================================================================

We have computed the single inclusive gluon spectrum in dilute--dense scattering at next-to-leading order in the light-cone wave function approach. The soft gluon modes of the projectile in the window $\Lambda<k^+<\vee$ are dressed by a unitary operator $\Omega$, whose coefficients we fixed to order $g^3$ from the LCPT wave functions of the valence vacuum and of a single soft gluon. The spectrum is the expectation value of the number of dressed gluons in the scattered state. At order $g^4$ it consists of virtual terms (one gluon in the final state), squared real two-gluon amplitudes and their interference terms. We evaluated all of them in transverse coordinate space.

Our main results are the following.
\begin{enumerate}
\item \emph{Rapidity factorization at mid rapidity.} At leading logarithmic accuracy the interval $\delta Y=\log(\vee/\Lambda)$ opened by the extra gluon is shared, Eq.~\eqref{eq:LLmain}. The part above the measured gluon, $\log(\vee/k^+)$, evolves the projectile charge correlator with BFKL. The part below it, $\log(k^+/\Lambda)$, evolves the target Wilson lines with JIMWLK. Real and virtual terms combine into a double commutator of emission operators. Its action on the LO spectrum is exactly $-H_{\rm JIMWLK}$, including the terms in which the soft gluon is emitted by the measured gluon. All terms with three and four valence charges cancel. At mid rapidity, where both logarithms are large, the LO formula must therefore be dressed by both evolutions. The coefficient of $\log(\vee/\Lambda)$ alone is JIMWLK$\otimes$LO, not the sum of the two evolution kernels.
\item \emph{The finite part.} We defined the NLO remainder by an exact subtraction of the two logarithms, keeping all dependence on the finite cutoff $\vee$, and gave it term by term (Sec.~\ref{sec:finite}). Power-like terms $\propto1/\Lambda$ arise only in the rescattering term G1-III and cancel there. This requires the relative sign of the two orderings of the one-gluon wave function that is fixed by the cancellation of the instantaneous double pole at $p^+=k^+$, Eq.~\eqref{eq:endpoint}.
\item \emph{Running coupling and DGLAP.} The gluon splitting function governs the self-energy before the target, the loop across the target, and the initial-state, interference and final-state collinear splittings (weights $1:-2:1$). The $\beta_0$ poles of the self-energy and of the loop across the target come with opposite signs. Their balance leads to a ratio of three running couplings, Eq.~\eqref{eq:rcstruct}. When the projectile correlation length exceeds $1/Q_s$, part of the $\beta_0$ logarithm is projectile DGLAP evolution rather than running coupling~\cite{Kovner:2023vsy}. The final-state collinear pole renormalizes the gluon fragmentation function.
\end{enumerate}

\paragraph*{Status of the calculation.} The following items remain to be completed before the finite part can be evaluated numerically.
\begin{enumerate}
\item The diagonal term G1-I$'$ (the measured gluon emitting and reabsorbing a soft gluon after the target) is known at LL, where it is required for JIMWLK, but its finite part has not been computed. It must also cancel the edge logarithm $-2K_f\,\delta Y_T$ of G1-III.
\item G1-V part~2b has to be completed in dimensional regularization, and its overall weight relative to G1-II has to be fixed. This settles the $\beta_0$ pole balance, Sec.~\ref{sec:RG}, and the CMW constant.
\item The adjoint index placement of the Wilson lines in G2-I(2), G2-IV and G3-III(2) has to be made consistent for a general $U$, so that the cancellation of the collinear singularity among the real rows can be verified.
\item The Group~I bookkeeping of the $\log(\vee/\Lambda)$ terms has to include the commutator $[\bar{\mf N}_2,\bar{\mf A}^{(1)}]$ of Eq.~\eqref{eq:VII}. It contributes only at LL.
\item In a row-by-row comparison with the leading-log wave function, the three-charge LL terms of G3-III part~1 do not agree with the exact result. This row should be rechecked; its finite part, Eq.~\eqref{eq:finGIIIIII1}, depends on it.
\item The overall sign of the $p^+>k^+$ non-instantaneous one-gluon amplitude, Eq.~\eqref{eq:nonInst1}, is fixed here by the endpoint condition and by the cancellation of $1/\Lambda$. A direct re-derivation of its momentum-space form would confirm it.
\end{enumerate}
%TODO: items 1-6 above are the open items of the calculation.

\paragraph*{Outlook.} With these items completed, the NLO spectrum can be evaluated for a Gaussian (McLerran--Venugopalan) projectile and a target evolved with JIMWLK or BK, with the BFKL-evolved projectile correlator. The rapidity factorization~\eqref{eq:LLmain}, with the exact finite part and the explicit dependence on $\vee$, provides the natural framework for a consistent NLO treatment of gluon production at mid rapidity. It avoids the ambiguities of the rapidity subtraction encountered in forward production~\cite{Stasto:2013cha,Ducloue:2016shw,Iancu:2016vyg}. The inclusion of quarks ($n_f\neq0$) affects only the self-energy and the corresponding splitting functions.

\begin{acknowledgments}
[Acknowledgments to be added.]
\end{acknowledgments}
